% !TeX program = lualatex
% !TeX encoding = utf8
% !TeX spellcheck = uk_UA
% !TEX indexprogram = upmendex
% !TeX root =../ClassicalElectrodynamics.tex

% chktex-file 3
% chktex-file 7
% chktex-file 12
% chktex-file 21
% chktex-file 25


%% --------------------------------------------------------
\section{Таблиця переходу між гаусовою системою одиниць та SI}\label{sec:SGStoSI}
%% --------------------------------------------------------

Для коректного читання та інтерпретації літератури з класичної електродинаміки необхідно вільно орієнтуватися в різних системах одиниць, зокрема в системі СІ та гаусовій системі (СГС). Це пов'язано з тим, що різні джерела використовують різні системи одиниць: фундаментальні теоретичні праці традиційно формулюються в СГС, тоді як у сучасній навчальній і прикладній літературі переважає система СІ. Тому необхідно впевнено володіти перевідними співвідношеннями між цими системами для основних фізичних величин.

Перехід між системами одиниць не зводиться до простого масштабування величин, а супроводжується зміною вигляду рівнянь. У системі~SI рівняння містять константи \(\varepsilon_0\) та \(\mu_0\), тоді як у гаусовій системі вони відсутні, і відповідні коефіцієнти з'являються у вигляді множників \(4 \pi \). Тому необхідно враховувати не лише перевідні співвідношення, але й зміну структури рівнянь.

\begin{table}[h!]\centering
	\caption{Перехід між гаусовою системою одиниць та SI}\label{tab:CGStoSI}
	\small
	\begin{tblr}{
			colspec={l X[c] X[c]},
			row{1} = {1.5em, c, azure2, fg=white,font=\bfseries\sffamily},
			row{odd[2]}={gray9},
			row{2-Z}={1.7em},
			cell{1}{4}={mode=text},
			hlines={azure3},
			vlines={azure3},
			vline{2-4}={1}{white},
			hline{2}={white}
		}
		Величина                       & Система СГС               & Система SI                                                   \\
		Швидкість світла               & \(c\)                     & \((\varepsilon_0\mu_0)^{-1/2}\)                              \\
		Напруженість електричного поля & \(\Efield_{\text{СГС}}\)  & \(\frac{1}{\sqrt{4\pi\varepsilon_0}}\,\Efield_{\text{СГС}}\) \\
		Електричний потенціал          & \(\phi_{\text{СГС}}\)     & \(\frac{1}{\sqrt{4\pi\varepsilon_0}}\,\phi_{\text{СГС}}\)    \\
		Магнітна індукція              & \(\Bfield_{\text{СГС}}\)  & \(\frac{1}{\sqrt{4\pi\mu_0}}\,\Bfield_{\text{СГС}}\)         \\
		Вектор-потенціал               & \(\vect{A}_{\text{СГС}}\) & \(\frac{1}{\sqrt{4\pi\mu_0}}\,\vect{A}_{\text{СГС}}\)        \\
		Густина заряду                 & \(\rho_{\text{СГС}}\)     & \(\sqrt{4\pi\varepsilon_0}\,\rho_{\text{СГС}}\)              \\
		Заряд                          & \(q_{\text{СГС}}\)        & \(\sqrt{4\pi\varepsilon_0}\,q_{\text{СГС}}\)                 \\
		Густина струму                 & \(\vect{j}_{\text{СГС}}\) & \(\sqrt{4\pi\varepsilon_0}\,\vect{j}_{\text{СГС}}\)          \\
	\end{tblr}
\end{table}


%% --------------------------------------------------------
\subsection*{Приклад: переписування рівнянь Максвела з гаусової системи (СГС) в SI}
%% --------------------------------------------------------

Розглянемо перехід від гаусової системи (СГС) до системи~SI на прикладі рівнянь Максвелла. Це дозволяє простежити зміну вигляду рівнянь при переході між системами одиниць.

У гаусовій системі рівняння Максвела для вакууму мають вигляд:
\begin{equation*}
	\begin{aligned}
		\Div\Efield_{\text{СГС}} & = 4\pi\rho_{\text{СГС}},                                                                             & \qquad
		\Rot\Efield_{\text{СГС}} & = -\frac{1}{c}\frac{\partial \Bfield_{\text{СГС}}}{\partial t},                                               \\[4pt]
		\Div\Bfield_{\text{СГС}} & = 0,                                                                                                 & \qquad
		\Rot\Bfield_{\text{СГС}} & = \frac{4\pi}{c}\vect{j}_{\text{СГС}} + \frac{1}{c}\frac{\partial \Efield_{\text{СГС}}}{\partial t}.
	\end{aligned}
\end{equation*}

З таблиці беремо перевідні співвідношення (позначимо величини в SI тим же символом без індексу):
\begin{align*}
	 & \Efield_{\text{СГС}} = \sqrt{4\pi\varepsilon_0}\,\Efield, \quad
	\Bfield_{\text{СГС}} = \sqrt{4\pi\mu_0}\,\Bfield,                   \\
	 & \rho_{\text{СГС}} = \frac{\rho}{\sqrt{4\pi\varepsilon_0}}, \quad
	\vect{j}_{\text{СГС}} = \frac{\vect{j}}{\sqrt{4\pi\varepsilon_0}},  \\
	 & c = \frac{1}{\sqrt{\varepsilon_0\mu_0}}.
\end{align*}

Підставимо в перше рівняння (теорема Гауса):
\begin{equation*}
	\Div\bigl( \sqrt{4\pi\varepsilon_0}\,\Efield \bigr) = 4\pi \cdot \frac{\rho}{\sqrt{4\pi\varepsilon_0}}
	\;\Longrightarrow\;
	\sqrt{4\pi\varepsilon_0}\,\Div\Efield = \sqrt{\frac{4\pi}{\varepsilon_0}}\,\rho
	\;\Longrightarrow\;
	\Div\Efield = \frac{\rho}{\varepsilon_0}.
\end{equation*}

Друге рівняння (закон електромагнітної індукції Фарадея):
\begin{equation*}
	\Rot\bigl( \sqrt{4\pi\varepsilon_0}\,\Efield \bigr) = -\frac{1}{c}\frac{\partial}{\partial t}\bigl( \sqrt{4\pi\mu_0}\,\Bfield \bigr).
\end{equation*}
\begin{equation*}
	\sqrt{4\pi\varepsilon_0}\,\Rot\Efield = -\frac{\sqrt{4\pi\mu_0}}{c}\,\frac{\partial \Bfield}{\partial t}.
\end{equation*}
Оскільки \(c = 1/\sqrt{\varepsilon_0\mu_0}\), то \(\sqrt{\mu_0}/c = \sqrt{\mu_0}\sqrt{\varepsilon_0\mu_0} = \mu_0\sqrt{\varepsilon_0}\). Тоді:
\begin{equation*}
	\sqrt{4\pi\varepsilon_0}\,\Rot\Efield = -\sqrt{4\pi\varepsilon_0}\,\mu_0\,\frac{\partial \Bfield}{\partial t}
	\;\Longrightarrow\;
	\Rot\Efield = -\frac{\partial \Bfield}{\partial t}.
\end{equation*}

Третє рівняння:
\begin{equation*}
	\Div\bigl( \sqrt{4\pi\mu_0}\,\Bfield \bigr) = 0 \;\Longrightarrow\; \Div\Bfield = 0.
\end{equation*}

\noindent Четверте рівняння:
\begin{equation*}
	\Rot\bigl( \sqrt{4\pi\mu_0}\,\Bfield \bigr) = \frac{4\pi}{c}\cdot\frac{\vect{j}}{\sqrt{4\pi\varepsilon_0}} + \frac{1}{c}\frac{\partial}{\partial t}\bigl( \sqrt{4\pi\varepsilon_0}\,\Efield \bigr).
\end{equation*}
\begin{equation*}
	\sqrt{4\pi\mu_0}\,\Rot\Bfield = \frac{4\pi}{c\sqrt{4\pi\varepsilon_0}}\,\vect{j} + \frac{\sqrt{4\pi\varepsilon_0}}{c}\,\frac{\partial \Efield}{\partial t}.
\end{equation*}
Використаємо \(c = 1/\sqrt{\varepsilon_0\mu_0}\):

\begin{equation*}
	\frac{4\pi}{c\sqrt{4\pi\varepsilon_0}} = 4\pi\sqrt{\varepsilon_0\mu_0}\cdot\frac{1}{\sqrt{4\pi\varepsilon_0}} = \frac{4\pi\sqrt{\varepsilon_0\mu_0}}{\sqrt{4\pi\varepsilon_0}} = \sqrt{4\pi}\sqrt{\mu_0} = \sqrt{4\pi\mu_0}.
\end{equation*}
А другий доданок: \(\frac{\sqrt{4\pi\varepsilon_0}}{c} = \sqrt{4\pi\varepsilon_0}\sqrt{\varepsilon_0\mu_0} = \sqrt{4\pi}\,\varepsilon_0\sqrt{\mu_0}\).

Отже:
\begin{equation*}
	\sqrt{4\pi\mu_0}\,\Rot\Bfield = \sqrt{4\pi\mu_0}\,\vect{j} + \sqrt{4\pi}\,\varepsilon_0\sqrt{\mu_0}\,\frac{\partial \Efield}{\partial t}.
\end{equation*}
Поділимо на \(\sqrt{4\pi\mu_0}\):
\begin{equation*}
	\Rot\Bfield = \vect{j} + \frac{\sqrt{4\pi}\,\varepsilon_0\sqrt{\mu_0}}{\sqrt{4\pi\mu_0}}\,\frac{\partial \Efield}{\partial t}
	= \vect{j} + \varepsilon_0\sqrt{\mu_0}\cdot\frac{1}{\sqrt{\mu_0}}\,\frac{\partial \Efield}{\partial t}
	= \vect{j} + \varepsilon_0\,\frac{\partial \Efield}{\partial t}.
\end{equation*}

Отримані рівняння Максвелла в системі SI для вакууму:

\begin{equation*}
	\begin{aligned}
		 & \Div\Efield = \frac{\rho}{\varepsilon_0}, \qquad
		 & \Rot\Efield = -\frac{\partial \Bfield}{\partial t},                                    \\[4pt]
		 & \Div\Bfield = 0, \qquad
		 & \Rot\Bfield = \mu_0\vect{j} + \mu_0\varepsilon_0\,\frac{\partial \Efield}{\partial t}.
	\end{aligned}
\end{equation*}


Слід зауважити, що навіть у вакуумі в системі SI часто вводять допоміжні вектори електричної індукції \(\Dfield\) та напруженості магнітного поля \(\Hfield\), які для вакууму визначаються як:
\begin{equation*}
	\Dfield = \varepsilon_0 \Efield, \qquad \Hfield = \frac{\Bfield}{\mu_0}.
\end{equation*}
Підставляння цих визначень дозволяє переписати рівняння Максвелла в симетричнішій формі, де коефіцієнти \(\varepsilon_0\) та \(\mu_0\) зникають з правих частин\footnote{У такому записі рівняння Максвелла в SI набувають формально однакового вигляду як для вакууму, так і для середовищ, хоча у випадку середовищ \(\rho\) та \(\vect{j}\) позначають лише вільні заряди та струми, а поля \(\Dfield\) та \(\Hfield\) включають внесок зв'язаних зарядів і струмів намагніченості.}:
\begin{equation*}
	\begin{aligned}
		 & \Div\Dfield = \rho, \qquad
		 & \Rot\Efield = -\frac{\partial \Bfield}{\partial t},           \\[4pt]
		 & \Div\Bfield = 0, \qquad
		 & \Rot\Hfield = \vect{j} + \frac{\partial \Dfield}{\partial t}.
	\end{aligned}
\end{equation*}
